% GENERATED FILE. Edit canonical lessons, then run npm run generate:books.
% canonical-source-hash: fnv1a64:38717844544206af
% canonical-lessons: ZH-C444-qianshui, ZH-C444-dao, ZH-W444-dao-island, ZH-C444-bandao, ZH-C444-dengta, ZH-C444-fangshaishuang

\chapter{An Island}
\label{ch:zh-an-island}
\hlchaptermodality{444}

\begin{chapteropening}
\textbf{By the end of this chapter:} \emph{I can talk about diving, an island, a peninsula, a lighthouse and sunscreen.}

Complete the last lesson of chapter 444: I can talk about diving, an island, a peninsula, a lighthouse and sunscreen.
\end{chapteropening}

\section[qiánshuǐ]{qiánshuǐ — to go diving; to dive under water}

\label{lesson:ZH-C444-qianshui}



\begin{quote}
Before the new one: say the Mandarin for to escort someone, then the Mandarin for to catch fish.
\end{quote}

\subsection*{You'll want to know: qiánshuǐ}
\emph{qiánshuǐ} — \textquotedblleft{}to go diving; to dive under water\textquotedblright{}.

\begin{quote}
\emph{xiàtiān wǒmen qù hǎibiān qiánshuǐ} — In summer we go diving at the seaside.
\end{quote}

Its characters come later; for now it is a word for the ear.

\begin{cousinweb}
\emph{qiánshuǐ} is built from \emph{qián} 'to go under' and \zh{水} 'water', both inherited from Middle Chinese, \emph{dzjem} and \emph{sywijX} (Guangyun, in Baxter's transcription), which glosses the first as water flowing hidden and 'to hide'.
\end{cousinweb}

\begin{grammarlens}[title={What you've built}]
One more word for getting around.
\end{grammarlens}

\subsection*{Your turn}
\begin{itemize}
\raggedright
  \item \emph{Say it:} \emph{qiánshuǐ}
  \item \emph{Say it:} \emph{qiánshuǐ}, once more
  \item \emph{Say it:} \emph{xiàtiān wǒmen qù hǎibiān qiánshuǐ} — In summer we go diving at the seaside.
  \item \emph{Recall:} say the Mandarin for to escort someone, then the Mandarin for to catch fish, then say \emph{qiánshuǐ} again
\end{itemize}

\begin{tcolorbox}[breakable,colback=teal!4,colframe=teal!35!black,title={Before you move on}]
What does qiánshuǐ mean? (\textquotedblleft{}To go diving; to dive under water\textquotedblright{}.) Say it once more.
\end{tcolorbox}
\section[\texorpdfstring{\zh{岛} (dǎo)}{岛 (dǎo)}]{\zh{岛} (dǎo) — island}

\label{lesson:ZH-C444-dao}



\begin{quote}
Before the new one: say the Mandarin for to catch fish, then the Mandarin for to dive under water.
\end{quote}

\subsection*{You'll want to know: \zh{岛}}
\textbf{\zh{岛}} — \emph{dǎo} — \textquotedblleft{}island\textquotedblright{}.

\begin{quote}
\textbf{\zh{我们坐船去那个岛}} \emph{wǒmen zuò chuán qù nàge dǎo} — We're taking a boat to that island.
\end{quote}

The new character, in pieces: \zh{岛} — seven strokes in two pieces: on top a short stroke falling to the left, a small bar that turns down and hooks, a dot inside it, and a long stroke that runs down, turns right, turns down again and hooks; tucked under it \zh{山} shān, three strokes. The next lesson puts it in your hand.

\begin{cousinweb}
\textbf{\zh{岛}} (\emph{dǎo}) is inherited from Middle Chinese \emph{tawX} (Guangyun, in Baxter's transcription), which glosses it, after an older dictionary, as the hills here and there in the sea where one can stop and rest.
\end{cousinweb}

\begin{grammarlens}[title={What you've built}]
One more word for getting around.
\end{grammarlens}

\subsection*{Your turn}
\begin{itemize}
\raggedright
  \item \emph{Say it:} \emph{dǎo}
  \item \emph{Say it:} \emph{dǎo}, once more
  \item \emph{Say it:} \emph{wǒmen zuò chuán qù nàge dǎo} — We're taking a boat to that island.
  \item \emph{Recall:} say the Mandarin for to catch fish, then the Mandarin for to dive under water, then say \emph{dǎo} again
\end{itemize}

\begin{tcolorbox}[breakable,colback=teal!4,colframe=teal!35!black,title={Before you move on}]
What does \zh{岛} mean? (\textquotedblleft{}Island\textquotedblright{}.) Say it once more.
\end{tcolorbox}
\section[\texorpdfstring{\zh{岛} (dǎo)}{岛 (dǎo)}]{\zh{岛} — a top piece of four strokes, then \zh{山}}

\label{lesson:ZH-W444-dao-island}



\begin{quote}
The word first, then the character inside it.

\begin{itemize}
\raggedright
  \item \emph{Recall:} say \emph{dǎo} — \textbf{R1}, one lesson back, before it decays
\end{itemize}

\begin{quote}
\zh{山} \zh{海}
\end{quote}

\emph{Read it:} each one aloud, with nothing said first — \emph{shān}, \emph{hǎi}
\end{quote}

\subsection*{Script — a top piece of four strokes, then \zh{山}}
\hlblockfigure{\detokenize{figures/ZH-W444-dao-island-filmstrip.pdf}}{How \zh{岛} is written, stroke by stroke}

\begin{quote}
\zh{岛}
\end{quote}

\textbf{Two pieces, top before bottom.} Count the pieces, not the strokes.

\textbf{The top piece, four strokes.} A short stroke falling to the left. A small bar that turns down and hooks. A dot inside it. Then a long stroke that runs down the left, turns right along the foot, turns down again and ends in a hook.

\textbf{The bottom piece, \zh{山} shān, three strokes}, tucked into the corner: the middle vertical first, then the left one, which turns right along the foot, then the right one.

4 + 3 = seven.

\subsection*{Your turn}
\begin{itemize}
\raggedright
  \item \emph{Write it:} \zh{岛}, all seven, with the model in front of you
  \item \emph{Cover:} write it again from memory, and count the pieces, not the strokes
\end{itemize}

\begin{tcolorbox}[breakable,colback=teal!4,colframe=teal!35!black,title={Before you move on}]
How many strokes in \zh{岛}? (\textbf{Seven}.) Which written character sits at the bottom? (\textbf{\zh{山}}.) Which stroke of \zh{山} comes first? (\textbf{The middle vertical}.)

Source: \href{https://raw.githubusercontent.com/chanind/hanzi-writer-data/68d10a4b21150cae5e1ebbd223eed289cf32d90c/data/\%E5\%B2\%9B.json}{Hanzi Writer Data \zh{岛}, pinned stroke paths and medians}.
\end{tcolorbox}
\section[\texorpdfstring{\zh{半岛} (bàndǎo)}{半岛 (bàndǎo)}]{\zh{半岛} (bàndǎo) — peninsula}

\label{lesson:ZH-C444-bandao}



\begin{quote}
Before the new one: say the Mandarin for to dive under water, then the Mandarin for an island.
\end{quote}

\subsection*{You'll want to know: \zh{半岛}}
\textbf{\zh{半岛}} — \emph{bàndǎo} — \textquotedblleft{}peninsula\textquotedblright{}.

\begin{quote}
\textbf{\zh{这个半岛有很多酒店}} \emph{zhège bàndǎo yǒu hěn duō jiǔdiàn} — There are lots of hotels on this peninsula.
\end{quote}

\begin{cousinweb}
\textbf{\zh{半岛}} (\emph{bàndǎo}) is built from \zh{半} 'half' and \zh{岛} 'an island', both inherited from Middle Chinese, \emph{panH} and \emph{tawX} (Guangyun, in Baxter's transcription).
\end{cousinweb}

\begin{grammarlens}[title={What you've built}]
One more word for getting around.
\end{grammarlens}

\subsection*{Your turn}
\begin{itemize}
\raggedright
  \item \emph{Say it:} \emph{bàndǎo}
  \item \emph{Say it:} \emph{bàndǎo}, once more
  \item \emph{Say it:} \emph{zhège bàndǎo yǒu hěn duō jiǔdiàn} — There are lots of hotels on this peninsula.
  \item \emph{Recall:} say the Mandarin for to dive under water, then the Mandarin for an island, then say \emph{bàndǎo} again
\end{itemize}

\begin{tcolorbox}[breakable,colback=teal!4,colframe=teal!35!black,title={Before you move on}]
What does \zh{半岛} mean? (\textquotedblleft{}Peninsula\textquotedblright{}.) Say it once more.
\end{tcolorbox}
\section[dēngtǎ]{dēngtǎ — lighthouse}

\label{lesson:ZH-C444-dengta}



\begin{quote}
Before the new one: say the Mandarin for an island, then the Mandarin for a peninsula.
\end{quote}

\subsection*{You'll want to know: dēngtǎ}
\emph{dēngtǎ} — \textquotedblleft{}lighthouse\textquotedblright{}.

\begin{quote}
\emph{wǎnshang chuán wǎng dēngtǎ kāi} — At night the boats steer for the lighthouse.
\end{quote}

Its characters come later; for now it is a word for the ear.

\begin{cousinweb}
\emph{dēngtǎ} is built from \emph{dēng} 'a lamp' and \emph{tǎ} 'a tower', both inherited from Middle Chinese, \emph{tong} and \emph{thap} (Guangyun, in Baxter's transcription), which glosses the first as a lamp and the second as a pagoda.
\end{cousinweb}

\begin{grammarlens}[title={What you've built}]
One more word for getting around.
\end{grammarlens}

\subsection*{Your turn}
\begin{itemize}
\raggedright
  \item \emph{Say it:} \emph{dēngtǎ}
  \item \emph{Say it:} \emph{dēngtǎ}, once more
  \item \emph{Say it:} \emph{wǎnshang chuán wǎng dēngtǎ kāi} — At night the boats steer for the lighthouse.
  \item \emph{Recall:} say the Mandarin for an island, then the Mandarin for a peninsula, then say \emph{dēngtǎ} again
\end{itemize}

\begin{tcolorbox}[breakable,colback=teal!4,colframe=teal!35!black,title={Before you move on}]
What does dēngtǎ mean? (\textquotedblleft{}Lighthouse\textquotedblright{}.) Say it once more.
\end{tcolorbox}
\section[fángshàishuāng]{fángshàishuāng — sunscreen}

\label{lesson:ZH-C444-fangshaishuang}



\begin{quote}
Before the new one: say the Mandarin for a peninsula, then the Mandarin for a lighthouse.
\end{quote}

\subsection*{You'll want to know: fángshàishuāng}
\emph{fángshàishuāng} — \textquotedblleft{}sunscreen\textquotedblright{}.

\begin{quote}
\emph{wǒ de fángshàishuāng bú jiàn le} — My sunscreen has gone missing.
\end{quote}

Its characters come later; for now it is a word for the ear.

\begin{cousinweb}
\emph{fángshàishuāng} is built from \emph{fáng} 'to guard against', \emph{shài} 'to put in the sun' and \emph{shuāng} 'frost', all inherited from Middle Chinese, \emph{bjang}, \emph{sreaH} and \emph{srjang} (Guangyun, in Baxter's transcription); the sense 'a cream' for the last is modern.
\end{cousinweb}

\begin{grammarlens}[title={What you've built}]
One more word for getting around.
\end{grammarlens}

\subsection*{Your turn}
\begin{itemize}
\raggedright
  \item \emph{Say it:} \emph{fángshàishuāng}
  \item \emph{Say it:} \emph{fángshàishuāng}, once more
  \item \emph{Say it:} \emph{wǒ de fángshàishuāng bú jiàn le} — My sunscreen has gone missing.
  \item \emph{Recall:} say the Mandarin for a peninsula, then the Mandarin for a lighthouse, then say \emph{fángshàishuāng} again
\end{itemize}

\begin{tcolorbox}[breakable,colback=teal!4,colframe=teal!35!black,title={Before you move on}]
What does fángshàishuāng mean? (\textquotedblleft{}Sunscreen\textquotedblright{}.) Say it once more.
\end{tcolorbox}
